% SOLUCIÓN del Ejercicio 9 (para el instructor)
\documentclass[stu,12pt,a4paper,floatsintext]{apa7}
\usepackage[utf8]{inputenc}
\usepackage[T1]{fontenc}
\usepackage[spanish,es-tabla]{babel}
\usepackage{csquotes}
\usepackage{booktabs}
\usepackage[style=apa,backend=biber,sortcites=true,sorting=nyt]{biblatex}
\DeclareLanguageMapping{spanish}{spanish-apa}
\addbibresource{referencias_apa.bib}

\title{Acondicionamiento de señales con amplificadores operacionales}
\shorttitle{Acondicionamiento de señales}
\authorsnames{Estudiante Ejemplo}
\authorsaffiliations{Universidad Nacional Tecnológica de Lima Sur}
\course{IE-205: Electrónica Analógica}
\professor{Docente Ejemplo}
\duedate{27 de octubre de 2026}
\abstract{Este informe breve resume cómo un amplificador operacional en configuración
inversora permite acondicionar señales de bajo nivel antes de su digitalización. Se
revisan la ganancia en lazo cerrado, el efecto del ancho de banda y los criterios para
elegir resistencias, y se contrastan los resultados con la bibliografía de referencia.
El trabajo también ilustra el uso de LaTeX con el estilo APA 7: portada, citas
narrativas y parentéticas, tablas sin líneas verticales y lista de referencias
generada automáticamente. Se declara además el uso de una herramienta de inteligencia
artificial para depurar el código, de acuerdo con la política de integridad académica.
Finalmente, se discuten limitaciones del modelo ideal y recomendaciones prácticas para
laboratorio. Los resultados muestran que la ganancia teórica se cumple dentro de la tolerancia de los componentes mientras la frecuencia de operación permanezca muy por debajo del producto ganancia-ancho de banda del dispositivo, lo que respalda el uso del modelo ideal como primera aproximación de diseño.}
\keywords{amplificador operacional, ganancia, acondicionamiento de señales}

\AtBeginDocument{\renewcommand{\notelabel}{Nota}}
\begin{document}
\maketitle

\section{Introducción}
El amplificador operacional es un bloque fundamental del análisis de circuitos
analógicos \parencite{sedra2020}. Como señala \textcite{boylestad2023}, la ganancia
en lazo cerrado depende de la razón entre resistencias.

El formato de este informe sigue el manual de la \textcite{apa2020}.

\section{Resultados}
La Tabla~\ref{tab:ganancia} muestra la ganancia para tres valores de $R_f$.

\begin{table}[htbp]
  \caption{Ganancia del amplificador inversor}
  \label{tab:ganancia}
  \centering
  \begin{threeparttable}
  \begin{tabular}{cc}
    \toprule
    $R_f$ (k$\Omega$) & Ganancia $-R_f/R_1$ \\
    \midrule
    10 & $-1$ \\
    47 & $-4{,}7$ \\
    100 & $-10$ \\
    \bottomrule
  \end{tabular}
  \begin{tablenotes}[para]
    \item \textit{Nota.} Con $R_1 = 10$~k$\Omega$.
  \end{tablenotes}
  \end{threeparttable}
\end{table}

\section{Uso de inteligencia artificial}
Se utilizó ChatGPT \parencite{openai_chatgpt} para depurar errores de compilación
del código \LaTeX. Las fuentes fueron verificadas por el autor.

\printbibliography[title={Referencias}]
\end{document}
